\section*{Resumen.}

Este proyecto presenta el diseño e implementación de un sistema web para la gestión integral de un taller de reparación de equipos de tecnologías de la información (TI). El objetivo principal del sistema es centralizar y automatizar procesos clave como la recepción de equipos, el diagnóstico técnico, la elaboración de cotizaciones, la autorización de trabajos, la reparación, el control de inventario, los pagos, la entrega y la gestión de garantías. La solución se desarrolla con Laravel y contempla módulos para usuarios, clientes, almacén, órdenes de servicio, reportes y control de accesos por roles.

La plataforma permite mejorar la trazabilidad de cada servicio, reducir errores en los registros, controlar de forma más eficiente refacciones y cobros, y ofrecer información actualizada en tiempo real a administradores, técnicos y clientes. El proyecto combina el análisis de requisitos, el diseño de la arquitectura de software, el modelado de datos y la implementación funcional, con el propósito de optimizar la operación del taller, elevar la calidad del servicio y fortalecer la seguridad y confiabilidad de la información.

